\subsection{Forensic Intervention Formulation}
\label{sec:formulation}

Our starting point is a distinction between \emph{predictive correlation} and
\emph{trace-specific evidence}. A pixel may receive a high manipulation score
because its residual is inconsistent with the rest of the image, but it may also
receive the same score because it belongs to a salient object boundary that is
over-represented in the training masks. Feature separation reduces the second
failure mode only indirectly. We instead ask whether a suspected trace is
necessary for the prediction. Let $\mathbf A=\{\mathbf A_i\}_{i=1}^L$ denote the
observed forensic pyramid and let $\mathcal I(\mathbf A,\mathbf m)$ replace traces
selected by proposal $\mathbf m$ without changing the RGB content path. We call
$\widetilde{\mathbf A}=\mathcal I(\mathbf A,\mathbf m)$ a trace counterfactual.
For a weight-shared verifier $V_\theta$, the pixel-wise effect is positive when
removing the trace withdraws manipulation evidence. This construction embodies
three controls: factual and counterfactual inputs come from the same image, their
content features are identical, and their verifier parameters are identical.

The proposal is necessarily imperfect. Hence, a single positive effect cannot be
treated as proof. \method{} maintains two hypotheses throughout decoding:
$H_1$ states that a local manipulation trace is present, whereas $H_0$ states that
the observed deviation is compatible with the pristine image. Every scale writes
support for both hypotheses. Their ratio yields a decision; their simultaneous
strength yields conflict. This causal-effect-and-ledger formulation determines
the order of all modules below: quality calibration must precede intervention,
intervention must precede evidence accounting, and ledger conflict must precede
uncertainty-guided refinement.

\textbf{Design principle and novelty boundary.}
The framework is deliberately not a collection of interchangeable attention
blocks. Its state changes have distinct meanings. QRC estimates whether a trace
observation is trustworthy; ICTI asks what the prediction would be without that
observation; the shared verifier measures the resulting withdrawal of evidence;
DHEL records support and refutation without cancelling either; and CBR spends
additional capacity where those accounts disagree. Removing an earlier stage
therefore changes the input semantics of every later stage. This ordering also
separates our objective from frequency selection, prototype matching, or generic
two-branch fusion: none of those operations by itself tests the necessity of a
suspected trace while preserving the content path and verifier. We use
``counterfactual'' in this operational, model-internal sense rather than claiming
to reconstruct the unknown unedited photograph. The intervention supplies a
controlled diagnostic contrast, and the ledger makes the limitations of that
contrast visible as conflict instead of converting it immediately into a hard
mask.

\subsection{Quality-Conditioned Trace Observation}
\label{sec:qrc}

\textbf{Exactly constrained residuals.}
The RGB path preserves scene structure, while the trace path begins with $M$
learnable $5\!\times\!5$ residual kernels. For the $m$th kernel and center
$p_0$, we impose
\begin{equation}
\sum_{c,p}\mathbf K_m(c,p)=0,\quad
\sum_c\mathbf K_m(c,p_0)=-1,\quad
\sum_{c,p\ne p_0}\mathbf K_m(c,p)=1.
\label{eq:new_constraint}
\end{equation}
The three center coefficients are $-1/3$; all non-center coefficients are a
softmax over unconstrained parameters. Equation~\eqref{eq:new_constraint} is
therefore satisfied on every forward pass rather than approximately restored
after optimization. The zero DC response suppresses locally constant appearance,
but it does not make residual strength comparable across unknown processing
chains.

\textbf{Calibration and evidential reliability.}
From grayscale image $g$, we measure mean, standard deviation, saturation ratio,
horizontal/vertical gradient energy, Laplacian energy, and nominal 8-pixel block
boundary energy:
\begin{equation}
\begin{split}
\mathbf q&=[\mu_g,\sigma_g,\rho_{\rm sat},E_x,E_y,E_{\rm lap},B_x,B_y],\\
\bar{\mathbf T}_i&=(1+0.5\tanh\bm\gamma_i)\odot\mathbf T_i
                  +0.5\tanh\bm\beta_i,\quad
\rho_i=\sigma(r_i),
\end{split}
\label{eq:new_quality}
\end{equation}
where $(\bm\gamma_i,\bm\beta_i,r_i)$ are produced by a small MLP from $\mathbf q$.
The FiLM terms calibrate the observation; scalar $\rho_i\in(0,1)$ states how much
the ledger should trust scale $i$. This division is important. A strongly
recompressed image can still contain a large high-frequency response, yet that
response should receive low evidential weight. The final MLP layer is initialized
to zero, making FiLM an identity and $\rho_i=0.5$ at the beginning of training.

\textbf{Content leakage suppression.}
QRC is followed by a lightweight safeguard rather than treated as a principal
contribution. Aligned trace and content vectors are locally projected:
\begin{equation}
\mathbf T_i^\perp=\bar{\mathbf T}_i-\sigma(a_i)
\frac{\langle\bar{\mathbf T}_i,\mathbf C_i\rangle}
{\|\mathbf C_i\|_2^2+\epsilon}\mathbf C_i,\qquad
\mathbf A_i=\operatorname{LN}(\mathbf W_i*\mathbf T_i^\perp).
\label{eq:new_orth_anchor}
\end{equation}
The projected and original traces are mixed before $\mathbf W_i$, so useful
content-coupled acquisition statistics are not irreversibly deleted. All anchors
have width $d$. We do not claim that orthogonality alone establishes forensic
causality; it merely reduces leakage before the explicit intervention test.

The distinction between calibration and reliability is also consequential for
optimization. If quality descriptors directly gated the final mask, the model
could learn dataset-specific correlations such as ``strong compression implies
tampering.'' Here they can only rescale trace observations and discount evidence
writes. A low $\rho_i$ cannot create pristine evidence, nor can it suppress the
RGB content path; it merely limits how strongly a questionable forensic scale
votes in DHEL. Conversely, FiLM cannot decide the final class because its output
must still survive intervention and shared verification. The two outputs of QRC
thus constrain one another: calibration improves observability, while reliability
prevents an amplified but unstable response from dominating the ledger. Per-scale
reliability further permits a blurred image to distrust the finest scale while
retaining useful block or region statistics at coarser scales.

\subsection{Intra-Image Counterfactual Trace Intervention}
\label{sec:icti}

ICTI constructs an alternative trace pyramid that is plausible under $H_0$ and
matched in semantic content. It contains a proposal, an image-internal pristine
estimator, and a soft intervention. No external generator or pristine reference
image is needed.

\textbf{Proposal and intervention curriculum.}
The finest anchor yields proposal logits $\mathbf a=P(\mathbf A_1)$ and candidate
probability $\mathbf m=\sigma(\mathbf a)$. During training, early proposals may
miss the entire manipulated region, in which case the counterfactual would equal
the factual input and provide no gradient. With probability $p_e=p_0\eta^e$ at
epoch $e$, we mix the resized ground-truth mask into the intervention selector:
\begin{equation}
\mathbf m_i^{(e)}=(1-\lambda_t)\mathcal R_i(\sigma(\mathbf a))
 +\lambda_t\mathcal R_i(\mathbf y),\qquad p_e=p_0\eta^e.
\label{eq:new_curriculum}
\end{equation}
We use $p_0=0.6$, $\eta=0.95$, and $\lambda_t=0.5$. When teacher forcing is not
sampled, the first term is used alone; inference never observes $\mathbf y$.
Proposal BCE and Dice supervision ensure that the selector remains meaningful
after the curriculum decays.

\textbf{Pristine trace estimation from the same image.}
Let $\mathbf b_i=1-\operatorname{sg}(\mathbf m_i)$ be detached background weights.
For pixel $p$, local and global pristine estimates are
\begin{equation}
\mathbf l_i(p)=
\frac{\sum_{q\in\mathcal N(p)}b_i(q)\mathbf A_i(q)}
     {\sum_{q\in\mathcal N(p)}b_i(q)+\epsilon},\qquad
\mathbf g_i=
\frac{\sum_q b_i(q)\mathbf A_i(q)}{\sum_q b_i(q)+\epsilon}.
\label{eq:new_pristine}
\end{equation}
The local estimate preserves spatially varying camera and compression statistics;
the global estimate is a stable fallback when $\mathcal N(p)$ is mostly
suspicious. Detaching $\mathbf b_i$ prevents the proposer from changing merely to
make pristine estimation easier. A gated convolution
$R_i=R([\mathbf A_i,\mathbf l_i,\mathbf g_i,\mathbf m_i])$ selects between the two
contexts and adapts them to the candidate location. The adapter is bounded to
$0.1\tanh(\cdot)$ around their convex mixture; it therefore corrects estimation
bias without becoming an unconstrained feature generator.

\textbf{Soft trace erasure.}
A learned gate $h_i\in(0,1)$ controls intervention strength. The counterfactual
anchor and observable intervention energy are
\begin{equation}
\widetilde{\mathbf A}_i=\mathbf A_i+\mathbf m_i h_i\odot
(\mathbf R_i-\mathbf A_i),\qquad
\mathbf e_i=\operatorname{mean}_c|\mathbf A_i-\widetilde{\mathbf A}_i|.
\label{eq:new_intervention}
\end{equation}
Hard replacement would create a new edge at the proposal boundary and allow the
detector to recognize its own intervention. The soft form makes replacement
continuous and exposes $\mathbf m_i h_i$ for a background-locality penalty. Since
only trace anchors are altered, both branches reuse the exact RGB content
features already absorbed during leakage suppression.

ICTI contains three safeguards against a trivial counterfactual. First, the
replacement is estimated only from the current image, so the branch cannot use
the identity of an external donor as an easy cue. Second, the selector is
supervised but detached when it defines pristine weights; otherwise it could
expand the presumed-pristine set solely to minimize reconstruction error. Third,
the residual adapter and intervention gate are bounded, while off-target energy
is penalized. An all-zero intervention is rejected by proposal supervision and
the positive-region ranking loss; an all-image intervention is rejected by
background stability and locality; an adversarial replacement is restricted by
the convex local/global estimate and bounded adapter. These controls do not prove
that $\widetilde{\mathbf A}_i$ is the true pristine feature. They make the learned
contrast falsifiable: a useful intervention must be localized, lower the factual
score within manipulated pixels, preserve background predictions, and remain
stable under benign degradation.

\subsection{Strictly Shared Causal Verification}
\label{sec:verifier}

The factual and counterfactual anchors are independently transported by the same
multi-scale operator. Starting from $\mathbf H_i^0=\mathbf A_i$ (or
$\widetilde{\mathbf A}_i$), each iteration re-injects its immutable anchor and
exchanges neighboring-scale context:
\begin{equation}
\begin{split}
\overleftarrow{\mathbf H}_i^k&=\Phi_{\rm td}(\mathbf H_i^{k-1},
\mathcal R(\overleftarrow{\mathbf H}_{i+1}^k),\mathbf H_i^0),\\
\mathbf H_i^k&=\Phi_{\rm bu}(\overleftarrow{\mathbf H}_i^k,
\mathcal R(\mathbf H_{i-1}^k),\mathbf H_i^0).
\end{split}
\label{eq:new_transport}
\end{equation}
$\Phi_{\rm td}$ and $\Phi_{\rm bu}$ are gated ConvNeXt-style cells shared across
scales and across the two branches. We use two iterations. The shared verifier
$V_\theta$ then produces factual and counterfactual logits. Their controlled
effect is
\begin{equation}
\Delta(p)=V_\theta(\{\mathbf H_i^K\})(p)-
V_\theta(\{\widetilde{\mathbf H}_i^K\})(p).
\label{eq:new_effect}
\end{equation}
Because the branches differ only through Eq.~\eqref{eq:new_intervention}, a
positive $\Delta$ is evidence that the erased trace supported the original
decision. Shared weights are essential: independent decoders could manufacture a
large $\Delta$ through parameter asymmetry. We inject only
$\operatorname{ReLU}(\Delta)$ into the final manipulation logit. Negative effects
are not discarded, however; they enter the refutation account below.

\subsection{Dual-Hypothesis Evidence Ledger}
\label{sec:dhel}

\textbf{Local evidence writes.}
At every scale, a shared writer compares the factual state, counterfactual state,
and their signed difference. Two softplus outputs create non-negative support
$s_i$ and refutation $r_i$:
\begin{equation}
[s_i,r_i]=\rho_i\operatorname{softplus}\!\left(
W_E*[\mathbf H_i^K,\widetilde{\mathbf H}_i^K,
\mathbf H_i^K-\widetilde{\mathbf H}_i^K]\right)+\epsilon.
\label{eq:new_write}
\end{equation}
The writer is not constrained to interpret every feature difference as support.
For example, a large difference that fails to lower the factual logit can be
written as refutation. Multiplication by $\rho_i$ connects acquisition quality to
evidential strength rather than to the final probability directly.

\textbf{Bounded bidirectional accounting.}
Let $\mathbf L_i=[s_i,r_i]$. Coarse-to-fine transport supplies region context;
fine-to-coarse transport returns boundary counter-evidence:
\begin{equation}
\overleftarrow{\mathbf L}_i=\mathbf L_i+\sigma(\alpha_i)
\mathcal R(\overleftarrow{\mathbf L}_{i+1}),\qquad
\mathbf L_i^*=\overleftarrow{\mathbf L}_i+\sigma(\beta_i)
\mathcal R(\mathbf L_{i-1}^*).
\label{eq:new_ledger_transport}
\end{equation}
The sigmoid-bounded discounts prevent one scale from contributing unbounded
evidence. Unlike feature addition, both ledger channels retain their identity
through transport. Consequently, a location with $(s,r)=(0.1,0.1)$ is distinct
from $(s,r)=(10,10)$ even though both have zero signed difference.

\textbf{Decision and conflict.}
For the aggregated support/refutation at each scale, we compute
\begin{equation}
\ell_i=\log(s_i^*+\epsilon)-\log(r_i^*+\epsilon),\qquad
\kappa_i=\frac{2\min(s_i^*,r_i^*)}{s_i^*+r_i^*+\epsilon}
\left(1-e^{-(s_i^*+r_i^*)}\right).
\label{eq:new_conflict}
\end{equation}
$\ell_i$ is a supervised manipulation log-odds; $\kappa_i\in[0,1]$ is high only
when both hypotheses have comparable and non-trivial strength. The exponential
mass factor prevents two negligible accounts from appearing maximally uncertain.
Five interpretable maps---normalized support, normalized refutation, conflict,
evidence mass, and intervention energy---are embedded and written back to the factual state. DHEL therefore retains an
auditable evidence representation while still providing rich decoder features.

This ledger differs from concatenating multi-scale features in two ways. First,
its channels have fixed hypothesis semantics at every resolution, so transport
cannot silently rotate support into refutation. Second, conflict depends on both
balance and evidence mass. Standard entropy is high for a probability near one
half even when the network has observed almost nothing; $\kappa_i$ instead becomes
large only when substantial evidence supports both hypotheses. This matters near
real splice boundaries, where an editing trace and a natural edge can coexist,
and in heavily compressed regions, where all accounts should remain weak. The
causal effect and intervention energy are included in the ledger writeback but do
not overwrite the accounts. DHEL can therefore express three different states:
confident manipulation, confident pristine content, and unresolved contradiction.
CBR consumes the third state explicitly rather than forcing DHEL to make a
premature binary decision.

\subsection{Conflict-Guided Boundary Refinement}
\label{sec:cbr}

The finest ledger-enriched state gives coarse logit $z^0$ and learned log-variance
$u^0$. Its probability gradient $b=\|\nabla\sigma(z^0)\|_2$ identifies a candidate
boundary. Learned uncertainty and ledger conflict are mixed, then supplied with
intervention energy to a residual head:
\begin{equation}
\begin{split}
q&=(1-\omega)\sigma(u^0)+\omega\kappa_1,
\quad \omega=\sigma(w),\\
(\delta z,\delta u)&=G([\mathbf H_1^{\rm led},\sigma(z^0),b,q,
\kappa_1,e_1]),\\
z&=z^0+(0.5+q)\delta z+\sigma(g)\operatorname{ReLU}(\Delta),
\quad u=\operatorname{logit}(q)+\delta u.
\end{split}
\label{eq:new_refine}
\end{equation}
This design makes uncertainty operational and partially mechanistic. High ledger
conflict expands the correction budget, whereas a strong RGB boundary with weak
intervention effect is not automatically amplified. The factual/counterfactual
effect is added only after their matched verification, preserving the controlled
contrast.

For image-level screening, reliability $c_p=\exp(-u_p)$ weights spatial evidence.
An MLP $h$ maps reliable mean and maximum to score $s$:
\begin{equation}
s=h\!\left(\frac{\sum_p c_p\sigma(z_p)}{\sum_p c_p+\epsilon},
\max_p c_p\sigma(z_p)\right).
\label{eq:new_image_score}
\end{equation}
Thus, ledger conflict affects both pixel refinement and the confidence assigned to
the whole-image decision.

\subsection{Non-Degenerate Causal Training}
\label{sec:loss}

\textbf{Localization and boundary terms.}
Small regions would otherwise be dominated by pristine pixels. With manipulated
area $A_+$ and background area $A_-$, define
$w_+=\operatorname{clip}(\sqrt{(A_-+1)/(A_++1)},1,w_{\max})$. The primary loss is
\begin{equation}
\mathcal L_{\rm het}=\frac1{HW}\sum_p w(y_p)e^{-u_p}
\operatorname{BCE}(z_p,y_p)+\tfrac12u_p,
\label{eq:new_hetero}
\end{equation}
combined with soft Dice and BCE between morphological gradients of prediction and
mask. Auxiliary decoder maps and proposal logits receive BCE+Dice supervision.

\textbf{Counterfactual validity and locality.}
The intervention must lower evidence inside a manipulated region yet leave the
background stable. For margin $\delta$, factual/counterfactual probabilities
$p^f,p^{cf}$, and logits $z^f,z^{cf}$, we use
\begin{equation}
\mathcal L_{\rm cf}=\frac{\sum_p y_p\operatorname{softplus}
(\delta-z_p^f+z_p^{cf})}{\sum_py_p+\epsilon}
+\frac{\sum_p(1-y_p)|p_p^f-p_p^{cf}|}{\sum_p(1-y_p)+\epsilon}.
\label{eq:new_cf_loss}
\end{equation}
A locality loss averages intervention strength outside the ground-truth region.
It prevents the counterfactual generator from altering the full feature map to
manufacture a large effect. Importantly, proposal supervision and locality are
both required: the former avoids an empty intervention; the latter avoids an
unbounded one.

\textbf{Ledger and transported causal consistency.}
Every ledger log-odds $\ell_i$ receives BCE+Dice supervision. For a differentiable
degradation $D$ that preserves geometry, we align both final masks and bounded
causal effects:
\begin{equation}
\begin{split}
\mathcal L_{\rm led}&=\frac1L\sum_i
[\operatorname{BCE}(\ell_i,\mathbf y)+\tfrac12\operatorname{Dice}(\ell_i,\mathbf y)],\\
\mathcal L_{\rm ct}&=\mathcal L_{\rm sym}(z(\mathbf x),z(D\mathbf x))
+\operatorname{SmoothL1}(\tanh\Delta(\mathbf x),
\tanh\Delta(D\mathbf x)).
\end{split}
\label{eq:new_ledger_consistency}
\end{equation}
Mask consistency alone can preserve a shortcut; effect consistency requires the
model to preserve which trace caused its decision after resizing, blur, or
quantization.

The objectives are optimized jointly because each closes a different escape
route. Pixel supervision alone permits semantic shortcut learning. Ranking alone
permits global or adversarial interventions. Locality alone permits an empty
intervention. Ledger supervision alone permits both accounts to grow without
testing their causal origin. Their combination requires a localized proposal, a
matched score decrease, stable background behavior, correctly signed evidence,
and degradation-stable attribution. Teacher assistance is annealed rather than
removed abruptly, allowing the proposal and verifier to co-adapt before inference
uses the proposal alone. We monitor manipulated-region causal drop,
counterfactual background change, intervention precision, and ledger
contradiction rate in addition to localization metrics; these diagnostics expose
degenerate solutions that a favorable F1 score could conceal.

The full objective is
\begin{equation}
\begin{split}
\mathcal L={}&\lambda_h\mathcal L_{\rm het}+\lambda_d\mathcal L_{\rm dice}
+\lambda_b\mathcal L_{\rm bd}+\lambda_p\mathcal L_{\rm prop}
+\lambda_{cf}\mathcal L_{\rm cf}\\
&+\lambda_{led}\mathcal L_{\rm led}+\lambda_{loc}\mathcal L_{\rm loc}
+\lambda_{ct}\mathcal L_{\rm ct}+\lambda_o\mathcal L_{\perp}
+\lambda_i\mathcal L_{\rm img}+\lambda_v\|\mathbf u\|_2^2.
\end{split}
\label{eq:new_total}
\end{equation}
The implementation exposes every coefficient and mechanism. Disabling ICTI also
disables counterfactual and locality losses; disabling DHEL removes ledger loss;
and disabling CBR removes variance regularization. This coupling keeps cumulative
ablations semantically valid rather than leaving an objective active after its
associated mechanism has been removed.
